% Part: first-order-logic
% Chapter: introduction
% Section: syntax

\documentclass[../../../include/open-logic-section]{subfiles}

\begin{document}

\olfileid{fol}{int}{syn}

\olsection{વાક્યરચના}

સંકેતોની કઈ શ્રેણીઓ પ્રથમ-ક્રમ તર્કશાસ્ત્રનાં !!{sentence}s ગણાય તે
પહેલાં ચોક્કસ કરવું પડશે. આ કામ આપણે પછી કરીશું; હાલ તો દાખલાઓ વડે
આગળ વધીએ. પ્રથમ-ક્રમ તર્કશાસ્ત્રના પાયાના રચનાઘટકો---તેનું
સંકેતભંડોળ---બે ભાગમાં વહેંચાય છે. પહેલા ભાગમાં એવા સંકેતો છે જે
ચોક્કસ વસ્તુઓ કહેવા અથવા ચોક્કસ વસ્તુઓને નિર્દેશવા વાપરીએ છીએ.
વસ્તુઓને નિર્દેશવા આપણે !!{constant}s વાપરીએ છીએ અને નિર્દેશેલી
વસ્તુઓ વિશે કંઈક કહેવા !!{predicate}s વાપરીએ છીએ. દાખલા તરીકે, એક
વસ્તુને નિર્દેશવા $\Obj a$ને !!a{constant} તરીકે વાપરી શકીએ અને પછી
!!{sentence}~$\Atom{\Obj P}{\Obj a}$ વડે તેના વિશે કંઈક કહી શકીએ.
``$\Obj a$'' અને ``$\Obj P$'' માટે તમારા મનમાં અર્થ હોય, તો તમે
$\Atom{\Obj P}{\Obj a}$ને અંગ્રેજી વાક્ય તરીકે વાંચી શકો (અને
રૂપલક્ષી તર્કશાસ્ત્ર પ્રથમ વાર શીખ્યા ત્યારે કદાચ એવું કર્યું પણ હશે).
પ્રથમ-ક્રમ તર્કશાસ્ત્રનાં આવાં સાદાં !!{sentence}s મળી જાય પછી
સંકેતભંડોળના બીજા ભાગ---તાર્કિક સંકેતો (સંયોજકો અને પરિમાણકો)---વડે
તેમાંથી વધુ જટિલ વાક્યો બનાવી શકીએ. તેથી, દાખલા તરીકે,
$(\Atom{\Obj P}{\Obj a} \land \Atom{\Obj Q}{\Obj b})$ અથવા
~$\lexists[\Obj x][\Atom{\Obj P}{\Obj x}]$ જેવાં પદો રચી શકીએ.

સખત અધિસિદ્ધાંતલક્ષી અભ્યાસ માટે જરૂરી અર્થવિચારની ચોક્કસ વ્યાખ્યાઓ
અને આપણા !!{derivation} તંત્રોના નિયમો આપી શકાય એ માટે સર્વપ્રથમ
પ્રથમ-ક્રમ તર્કશાસ્ત્રનું !!a{sentence} કોને ગણવું તેની ચોક્કસ વ્યાખ્યા
આપવી પડે. પાયાનો વિચાર સમજવો સહેલો છે: માત્ર !!{predicate}s અને
!!{constant}sમાંથી $\Atom{\Obj P}{\Obj a}$ જેવાં કેટલાંક સાદાં
!!{sentence}s રચી શકાય. પછી સંયોજકો અને પરિમાણકો વાપરીને તેમાંથી વધુ
જટિલ વાક્યો રચીએ. પરંતુ વધુ જટિલ !!{sentence}s રચવા કયા ચોક્કસ નિયમો
વાપરી શકીએ? આ નિયમો સ્પષ્ટ જણાવવા જ પડે, નહિ તો ``પ્રથમ-ક્રમ
તર્કશાસ્ત્રનું !!{sentence}'' આપણે પૂરતી ચોકસાઈથી વ્યાખ્યાયિત કર્યું
નથી. અહીં કેટલાક પ્રશ્નો છે. પહેલો પ્રશ્ન યોગ્ય સંકેતશ્રેણીઓને જ
!!{sentence}s ગણવાનો છે. બીજો પ્રશ્ન એવું કરવાની રીતનો છે કે
\emph{બધાં} !!{sentence}s વિશે ગાણિતિક સાબિતીઓ આપી શકાય. અંતે,
!!{sentence}s પરની ``~$!A$માં આવતા દરેક $\Obj x$ને $\Obj b$થી બદલો''
જેવી કેટલીક પાયાની ક્રિયાઓની પણ ચોક્કસ વ્યાખ્યાઓ આપવી પડશે. મુશ્કેલી એ
છે કે પ્રથમ-ક્રમ તર્કશાસ્ત્રમાં આવતા પરિમાણકો અને !!{variable}sને લીધે
આ કેવી રીતે કરવું તે સંપૂર્ણપણે દેખીતું નથી. દાખલા તરીકે,
$\lexists[\Obj x][\Atom{\Obj P}{\Obj a}]$ને !!a{sentence} ગણવું જોઈએ?
$\lexists[\Obj x][\lexists[\Obj x][\Atom{\Obj P}{\Obj x}]]$નું શું?
અને ``$(\Atom{\Obj P}{\Obj x} \land \lexists[\Obj x][\Atom{\Obj
P}{\Obj x}])$માં $\Obj x$ને $\Obj b$થી બદલો'' તેનું પરિણામ શું હોવું
જોઈએ?

\end{document}
